\section{Conceitos de Paralelismo}
Paralelismo e a tecnica de dividir uma computacao em partes que podem ser
executadas simultaneamente. Em vez de realizar todas as instrucoes de forma
sequencial em uma unica unidade de processamento, um programa paralelo tenta
explorar multiplas unidades de execucao para reduzir tempo de resposta, aumentar
vazao ou processar entradas maiores. Essa ideia se tornou central em computacao
de alto desempenho porque o aumento de frequencia dos processadores encontrou
limites praticos de consumo de energia, dissipacao de calor e complexidade de
projeto. Como observado por Sutter, a evolucao de desempenho passou a depender
cada vez menos de processadores sequenciais mais rapidos e cada vez mais do uso
efetivo de concorrencia e paralelismo \cite{sutterfree}. Essa mudanca tambem foi
discutida por Asanovic et al., que descrevem o paralelismo como um caminho
necessario para continuar aumentando desempenho em arquiteturas modernas
\cite{asanoviclandscape}.

Uma forma classica de classificar arquiteturas paralelas e a taxonomia de Flynn,
que organiza sistemas conforme o numero de fluxos de instrucoes e de dados
processados simultaneamente \cite{flynntaxonomy}. Nessa taxonomia, SISD
(\textit{Single Instruction, Single Data}) representa o modelo sequencial
tradicional, SIMD (\textit{Single Instruction, Multiple Data}) representa uma
mesma instrucao aplicada sobre varios dados, MISD (\textit{Multiple Instruction,
Single Data}) e uma categoria pouco comum na pratica, e MIMD (\textit{Multiple
Instruction, Multiple Data}) representa sistemas nos quais diferentes unidades
executam instrucoes distintas sobre dados distintos. CPUs multicore se aproximam
do modelo MIMD, enquanto GPUs modernas combinam elementos SIMD/MIMD com grande
quantidade de threads.

\section{Processamento Multicore}

Processadores multicore integram varios nucleos de CPU em um mesmo chip. Cada
nucleo e uma unidade de processamento relativamente complexa, capaz de executar
um fluxo de instrucoes com baixa latencia e bom desempenho em codigo geral. Em
comparacao com arquiteturas sequenciais antigas, um processador multicore permite
que varias threads ou processos avancem simultaneamente. Esse modelo se tornou a
estrategia dominante para continuar aumentando desempenho sem depender apenas do
aumento de frequencia \cite{hennessypatterson}.

Em programacao multicore, uma estrategia comum e criar threads que compartilham a
memoria do processo. Bibliotecas e modelos como POSIX Threads, OpenMP e C++
Threads permitem dividir lacos, tarefas ou regioes independentes do programa. O
desafio principal e garantir corretude e desempenho ao mesmo tempo. 

\section{Processamento Manycore}

Arquiteturas manycore aumentam de forma mais agressiva o numero de unidades de
execucao. Em vez de poucos nucleos complexos, elas utilizam muitos nucleos mais
simples, otimizados para alta vazao. O objetivo principal nao e minimizar a
latencia de uma unica thread, mas executar muitas threads em paralelo para
processar grandes volumes de dados. GPUs sao o exemplo mais comum de arquitetura
manycore, embora o conceito tambem apareca em outros aceleradores e coprocessadores.

A diferenca entre multicore e manycore nao e apenas quantitativa. Em uma CPU
multicore, cada nucleo tenta executar codigo geral de forma eficiente, mesmo
quando ha desvios, dependencias e acesso irregular a memoria. Em uma arquitetura
manycore, espera-se que o programa exponha uma grande quantidade de paralelismo.
Quando uma thread aguarda memoria, outras threads podem ser escalonadas para
ocupar a unidade de execucao. Dessa forma, o sistema esconde latencia por meio de
concorrencia massiva, em vez de depender apenas de caches grandes e controle
sofisticado \cite{kirkgpu}.

Esse modelo favorece problemas com alto paralelismo de dados e operacoes
semelhantes sobre muitos elementos. Simulacoes numericas, processamento de
imagens, algebra linear, aprendizado de maquina e certas etapas de processamento
de grafos podem se beneficiar desse tipo de arquitetura. Entretanto, problemas
irregulares exigem cuidado. Em grafos, por exemplo, vertices podem ter graus
muito diferentes, levando a desbalanceamento de carga. Alem disso, os acessos a
memoria dependem das arestas do grafo e podem ser pouco coalescidos, reduzindo a
eficiencia da memoria da GPU.

\section{GPUs}
GPUs (\textit{Graphics Processing Units}) surgiram como processadores
especializados para renderizacao grafica. Com o tempo, sua arquitetura evoluiu
para executar tambem computacao de proposito geral, area conhecida como GPGPU
(\textit{General-Purpose computing on Graphics Processing Units}). Essa evolucao
foi impulsionada pela alta largura de banda de memoria e pela capacidade de
executar milhares de threads leves em paralelo \cite{owensgpu, kirkgpu}.

No modelo CUDA da NVIDIA, a GPU executa kernels lancados pela CPU. Um kernel e
executado por uma grade de blocos; cada bloco contem varias threads. Threads de
um mesmo bloco podem cooperar por meio de memoria compartilhada e barreiras
locais. Internamente, as threads sao agrupadas em warps, que executam instrucoes
em estilo SIMD, mas especificamente SIMT (\textit{Single Instruction, Multiple Threads}). 
No SIMT, cada thread possui seu proprio contexto, mas threads de um mesmo warp executam a mesma
instrucao ao mesmo tempo quando seguem o mesmo fluxo de controle
\cite{cudaguide}.

Essa organizacao tem consequencias diretas para desempenho. Quando threads de um
warp seguem caminhos diferentes por causa de desvios condicionais, ocorre
divergencia, e os caminhos precisam ser serializados. Quando threads acessam
posicoes proximas e alinhadas da memoria global, a GPU consegue combinar acessos
em transacoes eficientes; quando os acessos sao dispersos, a largura de banda e
menos aproveitada. Portanto, algoritmos eficientes em GPU precisam expor muito
paralelismo, reduzir divergencia e organizar os acessos a memoria sempre que
possivel.

\subsection{GPUs de computacao}
Existem GPUs projetadas especificamente para cargas de computacao de alto desempenho, como
simulacoes, inteligencia artificial, analise de dados e processamente cientifico. Essas GPUs 
geralmente oferecem mais memoria, maior
largura de banda, suporte melhor a precisao dupla, recursos de virtualizacao,
correcao de erros de memoria e interconexoes de alta velocidade.

\leal{acho paia falar do rack se vai testar só em uma gb200, talvez mudar os numeros, será bom ter comparações 
para saber o significado real desses numeros}

Um exemplo recente desse tipo de sistema e o NVIDIA GB200 NVL72. Ele nao deve ser
entendido apenas como uma GPU isolada, mas como uma arquitetura em escala de rack.
Um unico rack integra 36 CPUs Grace e 72 GPUs Blackwell. A documentacao da
NVIDIA descreve o GB200 Grace Blackwell Superchip como a combinacao de uma CPU
Grace com duas GPUs Blackwell conectadas por NVLink-C2C. .
Dentro do rack NVL72, as GPUs sao interligadas por NVLink de
quinta geracao e por NVLink Switch, formando um dominio NVLink com 72 GPUs. A
ideia e que as GPUs do rack consigam se comunicar com baixa latencia e alta
largura de banda, aproximando o rack do comportamento de um acelerador
massivamente paralelo unico \cite{nvidiagb200nvl72}.

As larguras de banda mostram a diferenca entre os niveis da hierarquia. No nivel
de memoria local das GPUs, o GB200 NVL72 possui 13,4 TB de memoria HBM3E agregada
e 576 TB/s de largura de banda agregada de memoria. No nivel de comunicacao entre
GPUs, o dominio NVLink do rack oferece 130 TB/s de largura de banda agregada para
comunicacao GPU-GPU. Ja um superchip GB200, composto por uma CPU Grace e duas
GPUs Blackwell, oferece 372 GB de HBM3E, 16 TB/s de largura de banda de memoria e
3,6 TB/s de largura de banda NVLink. Assim, acessos a memoria local da GPU sao os
mais rapidos, comunicacoes entre GPUs dentro do mesmo dominio NVLink ainda sao de
altissima vazao, e comunicacoes para fora do rack passam a depender da rede do
cluster, como InfiniBand ou Ethernet de alto desempenho \cite{nvidiagb200nvl72}.
